% =====================================================================
% P5 merge snippet — M07: five-axis taxonomy of the field
% Source: manuscript_v2, Sec. 14 (sec:taxonomy, 3,731 words, 5 tables). This snippet is a
% CONDENSED version: the five axes, the dissimilarity, and how to use it — about one page.
% Full version available on request from the v2 source; not pasted here because it would
% duplicate \cref{sec:synthesis}.
% Adapted: \cref/\eqref, this paper's labels, their [nosep] enumerate style.
% NEW LABELS: sec:taxonomy, tab:taxonomy-axes, eq:taxdissim
% Suggested placement: a new section between \cref{sec:synthesis} and \cref{sec:open}.
% =====================================================================

\section{A taxonomy of the field, and when two studies are comparable}\label{sec:taxonomy}

The literature labels a study ``generalized thermoelastic'' or it does not. That single axis collapses
independent commitments into one word: a study may adopt a relaxation law for the heat flux, a gradient or
integral law for the stress, a coupling to a second field, and a validation regime, and each can be varied
while the others are held fixed. Two studies sharing the label can therefore differ in the mathematics of
their governing equations, in their modal content and in the evidence they offer. \Cref{sec:synthesis}
compares families; this section says what a family \emph{is} as a coordinate, so that a new model can be
placed without argument.

The heat-flux template of \cref{sec:foundations} is a special case of Axis~2: the operator is local in
space and first order in time, with the lags as the coordinate that separates the template's members. A
model that needs an integral, fractional or gradient operator has left the axis rather than changed a
parameter on it, which is why those models cannot be ranked against the template's members by setting a
lag to zero.

A study is classified by its vector of commitments $\mathbf{c}=(c_1,\dots,c_5)$, one coordinate per axis:

\begin{enumerate}[label=\textbf{Axis \arabic*:}]
\item \textbf{Physical axis (material class).} Classical elastic, thermoelastic, piezoelectric,
  piezothermoelastic, magneto-thermo-elastic, poroelastic, micropolar, two-temperature, semiconductor-like.
  It fixes which conserved quantities and driving fields exist before any approximation.
\item \textbf{Mathematical axis (operator structure).} Local PDE, integro-differential, fractional-order,
  nonlocal gradient, hereditary memory integral, or variational/weak form. This is the axis along which the
  most consequential differences run.
\item \textbf{Wave axis (modal content).} Does the thermal mode propagate at finite speed; how many coupled
  modes exist; is there attenuation; what is the dispersion character of each.
\item \textbf{Coupling axis (field topology).} Which fields are coupled and how strongly ---
  thermoelastic, thermo-piezoelectric, thermo-magneto-elastic, thermo-poro-elastic, thermo-micro-elastic.
  Coupling topology, not the choice of energy functional, decides whether an interface problem is tractable.
\item \textbf{Validation axis (evidence regime).} Existence and uniqueness only; closed-form solution for a
  special geometry; numerical solution validated against a reference; comparison with experiment; a shared
  benchmark suite.
\end{enumerate}

Two studies are \emph{directly} comparable when their vectors agree in the first two coordinates; if they
differ there, no comparison of numbers is meaningful and only a comparison of structure is. To make the
weaker statement quantitative, define the axis-weight dissimilarity
\begin{equation}\label{eq:taxdissim}
d(\mathbf{c},\mathbf{c}')=\sum_{j=1}^{5} w_j\,\bigl(1-\delta_{c_j c_j'}\bigr),
\qquad
w_{\mathrm{math}}\ge w_{\mathrm{phys}}>w_{\mathrm{coup}}>w_{\mathrm{wave}}>w_{\mathrm{val}} ,
\end{equation}
with the ordering stated because it is a modelling decision that can be contested: a difference in the
mathematical axis changes the governing operator and invalidates most direct comparisons, whereas a
difference in the validation axis changes only the epistemic status of the conclusion. The weights are
reported whenever the dissimilarity is used, and no result in this review depends on their numerical values.

\begin{table}[H]\centering\small
\caption{The five axes with their admissible values. A study occupies one cell of the product; the last
column gives an example assignment for a coupled fractional piezothermoelastic problem.}
\label{tab:taxonomy-axes}
\begin{tabularx}{\textwidth}{@{}p{1.6cm}Y Y@{}}
\toprule
\textbf{Axis} & \textbf{Admissible values} & \textbf{Example} \\
\midrule
Physical & elastic; thermoelastic; piezoelectric; piezothermoelastic; magneto-thermo-elastic; poroelastic;
  micropolar; two-temperature; semiconductor & piezothermoelastic \\
Mathematical & local PDE; integro-differential; fractional; nonlocal gradient; hereditary; variational &
  fractional + local coupling \\
Wave & finite/infinite thermal speed; number of coupled modes; attenuating/non-attenuating &
  two modes, one attenuated \\
Coupling & none; thermal; thermal + electric; thermal + magnetic; thermal + poro; thermal + micro-rotation &
  thermal + electric \\
Validation & proof only; closed form; verified numerics; experiment; shared benchmark &
  verified numerics \\
\bottomrule
\end{tabularx}
\end{table}

The taxonomy earns its place only if it changes a decision, and it does in three places in this review: it
is why \cref{sec:synthesis} compares families in the order it does; it is why \cref{sec:interfaces} insists
that an energy flux be defined consistently with the model before ratios are compared (see
\cref{sec:eqregister}); and it is why \cref{sec:open} separates questions that differ in the mathematical
axis from questions that differ only in evidence.
